\subsection{Map robustness and multi-receiver completion}\label{sec:map_revision}
\subsubsection{Incorrect relative mappings and shared renaming}
The 800-scene wrong-map audit freezes the joint scorers and changes 0, 2, 4, 8, 12, or 16 of 16 entries through bijections. Three interventions per budget are averaged within replication. At 25\% corruption, F1 falls from 0.7798 to 0.6422; at 50\%, it is 0.5474 (Figure~\ref{fig:review_mapping}a). All groups remain feasible, showing that constraints cannot repair class meaning.

The curve is nonmonotone because random permutations can retain relative agreement. A shared renaming changes every absolute entry but preserves all 800 partitions. Event F1 therefore cannot establish absolute naming accuracy.

\subsubsection{Independent adaptation and conditional map optimization}
Each replication supplies 40 adaptation scenes disjoint from all original splits and stress seeds, without refitting. One to three maps are hidden while receiver 0 stays known; unknown receivers receive zero to two anchors. Fixed prefix budgets 1, 5, and 40 are all reported, without test-selected stopping. The 27 configurations use ten seeds and reuse the initial 800 tests. This analysis followed inspection of those tests; the subsequent comparison uses independent tests.

For three unknown maps without anchors, one scene recovers 75/120 entries with F1 0.6787. Five scenes recover 120/120 and restore F1 0.7798 (Table~\ref{tab:review_mapping}), as do forty. Budgets two to four were untested, so five is only the first successful tested value. The ten map interventions do not constitute 120 independent Bernoulli trials or establish larger-vocabulary sample complexity.

Five-scene coordinate fitting recovers 60\% of entries with F1 0.7171, whereas exact fitting recovers all. In replication 0 with fixed groups, every single-map gain is nonpositive, but joint updating raises the objective from 201.0133 to 235.4834. Coupling between maps causes this trap despite complete source-domain enumeration. Exact sequential and alternating fitting tie in every tested configuration, with no added event-update gain.

\begin{table}[!t]\begin{adjustwidth}{-\extralength}{0cm}
\caption{Independent adaptation without target anchors. The reference map remains known. Counts sum ten interventions; F1 uses the same 800 tests. Coordinate updates one map at a time; sequential and alternating methods solve the exact conditional block. All receiver and gate violations are zero.}\label{tab:review_mapping}
\small\setlength{\tabcolsep}{3pt}\begin{tabularx}{\linewidth}{@{}rrrYrYr@{}}\toprule
Adaptation scenes & Unknown sources & Correct hidden entries & Sequential F1 & Feedback F1 [95\% CI] & Coordinate F1 & Rx / gate \\\midrule
1 & 1 & 35/40 & 0.7604 & 0.7604 [0.7280, 0.7929] & 0.7604 & 0 / 0 \\
1 & 2 & 61/80 & 0.7163 & 0.7163 [0.6628, 0.7697] & 0.6941 & 0 / 0 \\
1 & 3 & 75/120 & 0.6787 & 0.6787 [0.6272, 0.7301] & 0.6279 & 0 / 0 \\
5 & 1 & 40/40 & 0.7798 & 0.7798 [0.7713, 0.7883] & 0.7798 & 0 / 0 \\
5 & 2 & 80/80 & 0.7798 & 0.7798 [0.7713, 0.7883] & 0.7798 & 0 / 0 \\
5 & 3 & 120/120 & 0.7798 & 0.7798 [0.7713, 0.7883] & 0.7171 & 0 / 0 \\
40 & 1 & 40/40 & 0.7798 & 0.7798 [0.7713, 0.7883] & 0.7798 & 0 / 0 \\
40 & 2 & 80/80 & 0.7798 & 0.7798 [0.7713, 0.7883] & 0.7798 & 0 / 0 \\
40 & 3 & 120/120 & 0.7798 & 0.7798 [0.7713, 0.7883] & 0.7133 & 0 / 0 \\
\bottomrule\end{tabularx}

\end{adjustwidth}\end{table}

\widefig{review_mapping}{Mapping interventions. (a) Relative-map errors harm association, whereas shared renaming preserves it. (b,c) Missing fractions use all 16 entries, with the reference known and no target anchors. Five- and 40-scene curves coincide. F1 bars use replication means; hidden-entry fractions are descriptive. All violation rates are zero.}{fig:review_mapping}

Schema tests reject many-to-one maps, three-column coarse vectors, and a fifth unknown class with positive mass. This defines the implementation boundary, not open-set or hierarchical recovery. Supplementary Section S11 records anchors and maps.

\subsection{Independent comparison of missing-map estimators}\label{sec:deep_mapping}
Independent validation uses 400 adaptation and 800 test scenes, with ten replications of 40 and 80 scenes. Prefix budgets are 1, 5, and 40. Receiver 0 remains known; three unanchored four-class maps are hidden. All solvers share physical-only groups and frozen parameters. Reconstructing the boosted model from original training with its saved temperature reproduces all 800 partitions, with maximum probability difference below $4.5\times10^{-11}$.

Reference-only fitting uses receiver 0, while coordinates update one map at a time. Adapted spectral synchronization \cite{pachauri2013} selects pairwise relative maps, then uses a leading eigenspace and Hungarian rounding. Spectral-plus-coordinate refines the full RF objective. Exact fitting retains all pair-table scores; sequential fitting solves once, whereas alternation also proposes groups. All maps freeze before testing (Supplementary Section S15).

Five adaptation scenes give exact recovery of all 120 hidden entries for every scorer. Weighted, unweighted, and boosted F1 are 0.7786, 0.7997, and 0.8232, matching known-map references. These 120 entries belong to ten interventions, not 120 independent trials (Table~\ref{tab:deep_mapping}).

\begin{table}[!t]\begin{adjustwidth}{-\extralength}{0cm}
\caption{Missing-map estimators using five adaptation and 80 test scenes per replication. Three four-class maps are unknown; only the reference is known. Counts sum ten interventions; F1 intervals use ten replication means. Sequential exact, alternating exact, and known-map results coincide. All outputs satisfy receiver and gate constraints.}\label{tab:deep_mapping}
\small\setlength{\tabcolsep}{4pt}\begin{tabularx}{\linewidth}{@{}lYrr@{}}\toprule
Scorer & Map estimator & Correct / hidden & Event-pair F1 [95\% CI] \\\midrule
Weighted linear & Unresolved & 21/120 & 0.5316 [0.5119, 0.5514] \\
Weighted linear & Reference only & 113/120 & 0.7565 [0.7224, 0.7906] \\
Weighted linear & Spectral sync. & 114/120 & 0.7712 [0.7548, 0.7877] \\
Weighted linear & Spectral + coordinate & 114/120 & 0.7712 [0.7548, 0.7877] \\
Weighted linear & Coordinate & 70/120 & 0.7005 [0.6541, 0.7469] \\
Weighted linear & Exact block & 120/120 & 0.7786 [0.7718, 0.7853] \\
\addlinespace
Unweighted linear & Unresolved & 21/120 & 0.4695 [0.4452, 0.4937] \\
Unweighted linear & Reference only & 118/120 & 0.7912 [0.7713, 0.8110] \\
Unweighted linear & Spectral sync. & 120/120 & 0.7997 [0.7913, 0.8080] \\
Unweighted linear & Spectral + coordinate & 120/120 & 0.7997 [0.7913, 0.8080] \\
Unweighted linear & Coordinate & 66/120 & 0.7078 [0.6478, 0.7678] \\
Unweighted linear & Exact block & 120/120 & 0.7997 [0.7913, 0.8080] \\
\addlinespace
Boosted & Unresolved & 21/120 & 0.4458 [0.4230, 0.4686] \\
Boosted & Reference only & 114/120 & 0.7995 [0.7645, 0.8345] \\
Boosted & Spectral sync. & 114/120 & 0.8134 [0.7923, 0.8344] \\
Boosted & Spectral + coordinate & 114/120 & 0.8134 [0.7923, 0.8344] \\
Boosted & Coordinate & 69/120 & 0.7290 [0.6754, 0.7826] \\
Boosted & Exact block & 120/120 & 0.8232 [0.8183, 0.8282] \\
\addlinespace
\bottomrule\end{tabularx}

\end{adjustwidth}\end{table}

Exact-minus-spectral F1 is 0.0073 for weighted scoring with interval $[-0.0093,0.0239]$, and 0.0099 for boosting with $[-0.0125,0.0322]$. Unweighted results coincide, so no general superiority is established. At 40 scenes, reference-only, spectral, and exact fitting recover all maps, while coordinates can stall. Sequential and alternating maps coincide in all 90 scorer--budget--replication settings.

Two same-class conditions hide maps using five adaptation and 40 test scenes per replication. Exact fitting recovers every entry, yet at copy probability 0.65 the scorers give F1 0.4450, 0.4504, and 0.4537, equal to known-map references. Correct registration cannot distinguish nearby same-family events; the next test supplies all maps to isolate this ambiguity.

A five-scene weighted case has 160 grouped pairs and 13,824 states. Cached enumeration takes 2.44 ms versus 1413.08 ms for full reevaluation, with identical optima in three-run medians. A six-class, two-unknown kernel has 518,400 states, taking 377.79 ms for tables and 1.76 ms for search. These timings do not establish larger-vocabulary recovery; Supplementary Section S15 gives memory and allocation limits.

\subsection{Class evidence under harder descriptor conditions}\label{sec:deep_semantic}
Eight conditions use 400 new scenes each, with ten replications of 40 scenes. Four frozen scorers receive correct maps to isolate class-evidence quality. Semantic interventions preserve physical reports; dense conditions share a separate matched physical reference. Supplementary Section S15 specifies probability models, seeds, and intervals.

Nominal near-collisions have independent classes, with agreement probability $1/4$: 312 of 1303 copied-event parent pairs match in the new sample. Forcing copies to share the parent's class lowers weighted, unweighted, and boosted F1 from 0.7747, 0.7925, and 0.8226 to 0.7010, 0.7096, and 0.7279, respectively. Physical-only F1 remains 0.6773, attributing part of the semantic gain to class independence.

Increasing copy probability from 0.25 to 0.65 gives boosted F1 of 0.6335 with independent classes and 0.4482 with shared classes, versus physical-only F1 of 0.4460. The paired boosted-minus-physical gain is 0.0021 with interval $[-0.0113,0.0155]$; both linear intervals also include zero. Correct registration offers little event discrimination here.

\begin{table}[!t]\begin{adjustwidth}{-\extralength}{0cm}
\caption{Semantic stress tests: mean event-pair F1 across ten replication means, with 400 scenes per row and known maps. Semantic interventions preserve physical observations; dense rows share a separate physical reference. Error marginals match in law, though realized counts can differ. All outputs satisfy receiver and gate constraints. Supplementary Section S15 gives full intervals.}\label{tab:deep_semantic}
\small\setlength{\tabcolsep}{4pt}\begin{tabularx}{\linewidth}{@{}Yrrrr@{}}\toprule
Condition & Physical & Weighted linear & Unweighted linear & Boosted \\\midrule
Nominal independent classes & 0.6773 & 0.7747 & 0.7925 & 0.8226 \\
Same-class near collisions & 0.6773 & 0.7010 & 0.7096 & 0.7279 \\
Dense near collisions / independent classes & 0.4460 & 0.5831 & 0.6210 & 0.6335 \\
Dense near collisions / same class & 0.4460 & 0.4408 & 0.4471 & 0.4482 \\
Class prior 0.7/0.1/0.1/0.1 & 0.6773 & 0.7451 & 0.7606 & 0.7852 \\
Receiver-specific confusion and confidence & 0.6773 & 0.7601 & 0.7822 & 0.7992 \\
Independent errors / matched marginals & 0.6773 & 0.7714 & 0.7905 & 0.8184 \\
Event-shared errors / matched marginals & 0.6773 & 0.8334 & 0.8826 & 0.9081 \\
\bottomrule\end{tabularx}

\end{adjustwidth}\end{table}

Class imbalance and receiver-specific classifiers reduce gains (Table~\ref{tab:deep_semantic}). The latter retain error probability 0.12 but vary confusion destinations and confidence ranges, separating correct column registration from reliable recognition.

Sharing one error flag and wrong family across an event preserves marginal laws from an independent-error control. Weighted F1 rises from 0.7714 to 0.8334, a paired gain of 0.0620 with interval $[0.0532,0.0708]$. Agreement on a wrong family can preserve event compatibility, so event F1 does not certify class correctness or establish a universal effect of correlated errors.


\subsection{Risk under departures from the observation assumptions}\label{sec:risks_revision}
The first 20 tests per replication support nine interventions, two frozen scorers, and three decoders, totaling 10,800 outcomes. Clock equicorrelations 0.8 and $-0.25$ compare diagonal and matched declarations. Their small effects here do not establish robustness to arbitrary dependence.

Adding exponential delay with mean $8\,\mu$s to 30\% of genuine reports lowers gate recall from 0.9997 to 0.9353 and constrained F1 from 0.7822 to 0.7122. This is descriptor contamination, not measured NLOS validation.

The alarm test adds 12 false reports by copying three templates to all four receivers. Templates share time, carrier, bandwidth, duration, and canonical class vectors before independent Gaussian timing jitter of s.d. $0.2\,\mu$s and frequency jitter of 120 Hz. Native classes are re-encoded. The independent control uses separate templates with matched counts and marginal pool. Each false report has singleton truth; no common interfering waveform is generated.

Correlated false reports have clutter attachment 0.7821 versus 0.7061 in the independent control, yet F1 is also higher at 0.6023 versus 0.5105 (Table~\ref{tab:review_risks}). Both produce feasible false groups; dependence does not uniformly reduce F1.

With probability 0.2, a genuine report receives a same-receiver copy carrying its transmission label. The original stays unchanged; the copy gets independent Gaussian timing jitter of s.d. $0.15\,\mu$s and frequency jitter of 40 Hz, while bandwidth, duration, and class probabilities remain copied. These are detection records, not simulated waveform segments. Constrained F1 is 0.6009.

Let $n_{es}$ count reports of event $e$ at receiver $s$. The scene-level pair-recall ceiling under receiver exclusivity alone is
\[
 R_{\max}=\frac{\sum_e\sum_{s<t}\min(n_{es},n_{et})}
 {\sum_e\binom{\sum_s n_{es}}2}.
\]
The denominator includes same-receiver fragment pairs, which exclusive groups cannot recover. Each receiver pair permits at most $\min(n_{es},n_{et})$ links, attainable by layered grouping without gates. Empty-positive scenes receive ceiling zero. Figure~\ref{fig:review_risks}b averages scene ceilings within replication and then across ten replications. The fragment value 0.6529 bounds recall, not F1 (Supplementary Section S12).

\begin{table}[!t]\begin{adjustwidth}{-\extralength}{0cm}
\caption{Observation-assumption tests with ten replications of 20 scenes. Rx/gate entries are violation percentages; constrained zeros follow by construction. Clutter attachment counts false reports in non-singletons. The recall ceiling ignores gates and scores.}\label{tab:review_risks}
\small\setlength{\tabcolsep}{3pt}\begin{tabularx}{\linewidth}{@{}YYrrrr@{}}\toprule
Intervention & Event-pair F1 [95\% CI] & Gate recall & Clutter attachment & Recall ceiling & Rx / gate \\\midrule
Nominal subset & 0.7822 [0.7635, 0.8009] & 0.9997 & 0.0067 & 1.0000 & 0 / 0 \\
Positive corr.; diagonal declaration & 0.7816 [0.7618, 0.8014] & 0.9997 & 0.0067 & 1.0000 & 0 / 0 \\
Positive corr.; matched covariance & 0.7823 [0.7627, 0.8019] & 0.9997 & 0.0067 & 1.0000 & 0 / 0 \\
Negative corr.; diagonal declaration & 0.7796 [0.7612, 0.7979] & 0.9997 & 0.0057 & 1.0000 & 0 / 0 \\
Negative corr.; matched covariance & 0.7798 [0.7609, 0.7986] & 0.9997 & 0.0057 & 1.0000 & 0 / 0 \\
Extra positive delay & 0.7122 [0.6946, 0.7298] & 0.9353 & 0.0056 & 1.0000 & 0 / 0 \\
Independent false alarms; same count & 0.5105 [0.4993, 0.5217] & 0.9997 & 0.7061 & 1.0000 & 0 / 0 \\
Correlated false alarms & 0.6023 [0.5868, 0.6178] & 0.9997 & 0.7821 & 1.0000 & 0 / 0 \\
Fragmented reports & 0.6009 [0.5857, 0.6161] & 0.9998 & 0.0084 & 0.6529 & 0 / 0 \\
\bottomrule\end{tabularx}

\end{adjustwidth}\end{table}
\widefig{review_risks}{Assumption failures under frozen evidence. (a) Event F1 with 95\% replication intervals. (b) Gate recall and the exclusivity-only recall ceiling. Fragments can be unrepresentable despite admission, and correlated false reports can be feasible. Alarm controls have equal added-report counts.}{fig:review_risks}

\subsection{Medium-size optimization audit}\label{sec:medium_exact}
Report count alone selects the first two tests with at least 40 reports per replication. Nested random subsets of 20, 30, and 40 yield 60 MILPs under the frozen greedy objective. All 60 prove optimality within five-second limits at requested gap zero. Records retain incumbents, bounds, status, nodes, feasibility, and construction-plus-solve time.

Greedy mean objective losses are 0.18\%, 0.21\%, and 0.46\%, without implying a general bound. At 40 reports, restarts reduce loss to 0.11\%; greedy, restart, and exact F1 are 0.7974, 0.8098, and 0.8142 (Table~\ref{tab:review_exact}). Better objective search helps this sample but does not establish event truth.

\begin{table}[!t]\begin{adjustwidth}{-\extralength}{0cm}
\caption{Medium-size audit against proven optima. Times are medians in the review environment, including MILP construction. Every decoder is feasible; unproven incumbents are not labeled exact.}\label{tab:review_exact}
\small\setlength{\tabcolsep}{3pt}\begin{tabularx}{\linewidth}{@{}rrrrrrrY@{}}\toprule
$N$ & Proven optima & Greedy gap & Restart gap & Greedy F1 & Restart F1 & MILP F1 & Median ms: greedy / restart / MILP \\\midrule
20 & 20/20 & 0.2\% & 0.0\% & 0.7689 & 0.7750 & 0.7750 & 0.11 / 0.53 / 5.07 \\
30 & 20/20 & 0.2\% & 0.0\% & 0.7621 & 0.7649 & 0.7649 & 0.14 / 0.94 / 14.63 \\
40 & 20/20 & 0.5\% & 0.1\% & 0.7974 & 0.8098 & 0.8142 & 0.22 / 1.52 / 35.32 \\
\bottomrule\end{tabularx}

\end{adjustwidth}\end{table}

\paragraph{Fixed-registry processing time.}
Fresh report-path timing includes feature construction, scoring, and decoding for ten scenes at each of 40, 80, and 120 reports. Greedy medians are 0.64, 1.80, and 3.63 ms; triangle MILP medians are 34.56, 284.37, and 956.65 ms. All 30 solves prove optimality. The one-second cap excludes construction, which raises the maximum path to 1014.16 ms. These implementation timings assume no deployment deadline and exclude waveforms and offline map fitting. At 40 reports, greedy F1 is 0.8068 versus MILP F1 of 0.7781. Supplementary Section S15 records timings, solver status, and memory scope.

\subsection{Auxiliary public-recording metadata check}\label{sec:public_data}
\subsubsection{Source, inclusion rules, and independent capture blocks}
We use the publisher-selected six-file urban LoRaWAN sample \cite{loed2020,loeddataset}, released under CC BY 4.0 with recorded publisher MD5 and download SHA-256. Training uses 3 March 2019, 2 April 2019, and 2 May 2020; calibration uses 3 July 2020; testing uses 3 August and 2 September 2020. Recorded timestamps, including those preceding filename dates, are retained.

SHA-256 of the complete decoded payload defines the operational packet reference for labels and supervised fitting. Eligible records are CRC-valid unconfirmed uplinks with valid base64; message type comes from the payload header, avoiding the CSV's decimal-rendered binary labels. After exact deduplication, we exclude groups with repeated gateways, spans above 5 s, or identities shared across splits.

These rules remove obvious repetitions without independently proving one transmission. We therefore evaluate \emph{reference-packet association}, with split totals in Table~\ref{tab:loed_inclusion} and per-file exclusions in Supplementary Section S7.

\begin{table}[!t]\begin{adjustwidth}{-\extralength}{0cm}
\caption{LoED sample and inclusion counts. A visible reference packet has reports at two or more gateways. Eligible test reports come from two gateways. The original dataset covers nine gateways. This experiment does not use all nine.}\label{tab:loed_inclusion}
\small\renewcommand{\arraystretch}{1.12}\begin{tabular*}{\linewidth}{@{\extracolsep{\fill}}lrrrr@{}}\toprule
Split & Files & Raw reports & Eligible reports & Visible reference packets \\\midrule
Training & 3 & 260,231 & 39,447 & 257 \\
Calibration & 1 & 25,659 & 17,252 & 434 \\
Test & 2 & 40,980 & 23,677 & 1,417 \\
\bottomrule\end{tabular*}

\end{adjustwidth}\end{table}

Repeated payloads at individual gateways explain many 2019 exclusions. Results therefore concern identifiable receiver-exclusive references, not all traffic; CRC-invalid reports have no independent clutter labels. Reference hashes are stored separately, and scoring excludes addresses, frame counters, payloads and hashes, ports, and row indices.

\subsubsection{Measured-input adaptation and frozen parameters}
LoED has no waveform-derived uncertainty or class posteriors, so the reduced model uses reception time $t_i$, channel $F_i$, and decoded spreading factor $L_i$. Clock corrections $\widehat b_s$ use matching training references.

Gateway pairs need ten references for clock fitting. Median differences define a weighted least-squares system whose minimum-norm solution fixes the common time origin; unsupported corrections remain zero. The supported pair median is 0.999985 s, without test-identity fitting.

Corrected time is $t_i'=t_i-\widehat b_{s_i}$, with tolerance $G=\max(1\,\mathrm{ms},Q_{0.99})=3.02517\,\mathrm{ms}$. Here $Q_{0.99}$ is the 99th percentile of absolute calibration-reference residuals. This empirical tolerance does not carry the Gaussian guarantee of Proposition~\ref{prop:gate}. The features are
\begin{equation}
 \psi_{ij}=\left[\log\!\left(1+\frac{|t_i'-t_j'|}{G}\right),\;
 \ind[F_i\ne F_j],\;\ind[L_i\ne L_j]\right].\label{eq:loed_features}
\end{equation}
Timing scoring uses the first coordinate; metadata scoring uses all three. Inputs are standardized and logistic coefficients fitted on training data, with temperature fitted on separate calibration data. Admission requires $|t_i'-t_j'|\leq G$.

Reported channel frequency is not oscillator offset, and decoded spreading factor is not a probabilistic recognition output. We synthesize no geometry, RF uncertainty, or class posterior.

Corrected timestamps define fixed 60 s windows without identity-based reassignment. Each scorer uses connected components, pairwise assignment followed by union, and constrained decoding.

Evaluation retains singletons and counts window-split reference pairs as false negatives over each complete capture. We report captures separately and pool TP, FP, and FN. Two captures do not support deployment-population intervals based on individual pairs.

\subsubsection{Measured outcomes and interpretation}
The 23,677 reports contain 22,260 reference packets, including 1417 observed at both gateways. Metadata constraints produce 1411 true-positive, one false-positive, and six false-negative pairs: precision 0.9993, recall 0.9958, and F1 0.9975. On 3 August, TP is 900, FP 0, and FN 3 (F1 0.9983); on 2 September, TP is 511, FP 1, and FN 3 (F1 0.9961). Assignment and components give the same partition.

Timing-only components produce seven false positives and one receiver conflict. Assignment and constraints each produce five false positives with F1 0.9961 (Table~\ref{tab:loed_results}).

\begin{table}[!t]\begin{adjustwidth}{-\extralength}{0cm}
\caption{Public-recording results. Values are reference-packet F1. Rx counts clusters with repeated receivers. Within each scorer, all methods use the same observations, scores, gates, and capture windows. We pool pair counts before calculating F1. We do not average rounded scores.}\label{tab:loed_results}
\small\renewcommand{\arraystretch}{1.12}\begin{tabular*}{\linewidth}{@{\extracolsep{\fill}}llrrrr@{}}\toprule
Evidence & Decoder & 3 Aug. F1 & 2 Sep. F1 & Pooled F1 & Rx \\\midrule
Timing & Components & 0.9972 & 0.9922 & 0.9954 & 1 \\
Timing & Assignment & 0.9983 & 0.9922 & 0.9961 & 0 \\
Timing & Constrained & 0.9983 & 0.9922 & 0.9961 & 0 \\
Metadata & Components & 0.9983 & 0.9961 & 0.9975 & 0 \\
Metadata & Assignment & 0.9983 & 0.9961 & 0.9975 & 0 \\
Metadata & Constrained & 0.9983 & 0.9961 & 0.9975 & 0 \\
\bottomrule\end{tabular*}

\end{adjustwidth}\end{table}

Shared decoded meanings support a reproducible metadata-to-packet procedure, but metadata already resolves nearly all ambiguity. The two-gateway check has no measured class posteriors or source-specific protocol vocabulary and offers little room to distinguish decoder accuracy.

Assignment is already receiver-exclusive with two gateways, leaving higher-order grouping untested. Metadata F1 0.9975 does not validate multi-protocol robustness or RF uncertainty. Section~\ref{sec:public_sensitivity} tests the supervised preprocessing behind this result.

\subsection{Sensitivity of the auxiliary public check}\label{sec:public_sensitivity}
The public sensitivity audit fixes all 23,677 test IDs while varying training, calibration, clock correction, or quantile (Table~\ref{tab:review_public}). Original F1 is 0.9975, unchanged with May 2020 training alone. The two 2019 blocks provide one positive pair and no clock pair meeting ten-reference support, giving zero F1. Two settings lack usable training classes or calibration matches and are reported NE.

Without clock correction, metadata F1 is 0.9457. Quantile 0.95 gives F1 0.9933; quantile 0.999 expands the gate to 3.31 s and gives 0.9905. More candidates need not improve association. The original 0.99 rule stays fixed without test-F1 selection.

\begin{table}[!t]\begin{adjustwidth}{-\extralength}{0cm}
\caption{Public preprocessing sensitivity with unchanged test reports. Train positives count reference pairs available to scorer fitting. Clock corrections require at least ten matches for a receiver pair. $G$ is the empirical time gate. F1 is reference-packet F1. NE means the required fit cannot be estimated.}\label{tab:review_public}
\small\setlength{\tabcolsep}{3pt}\begin{tabularx}{\linewidth}{@{}Yrrrrr@{}}\toprule
Preprocessing setting & Train positives & $G$ (ms) & Aug. F1 & Sep. F1 & Pooled F1 \\\midrule
Original 3/1 split & 257 & 3.025 & 0.9983 & 0.9961 & 0.9975 \\
2019 training blocks only & 1 & 1003.010 & 0.0000 & 0.0000 & 0.0000 \\
May training only & 256 & 3.025 & 0.9983 & 0.9961 & 0.9975 \\
April training only & 0 & -- & NE & NE & NE \\
March train / April calibrate & -- & -- & NE & NE & NE \\
2019 train / May calibrate & 1 & 2000.887 & 0.0000 & 0.0000 & 0.0000 \\
Clock correction disabled & 257 & 1003.010 & 0.9603 & 0.9202 & 0.9457 \\
Calibration quantile 0.95 & 257 & 1.000 & 0.9956 & 0.9892 & 0.9933 \\
Calibration quantile 0.999 & 257 & 3307.621 & 0.9923 & 0.9875 & 0.9905 \\
\bottomrule\end{tabularx}

\end{adjustwidth}\end{table}

\subsection{Public-reference tails and clock supervision}\label{sec:deep_loed}
Of July's 434 reference pairs, four have corrected differences above two seconds despite matching payloads and counters; their reported channels differ. Five May pairs and two September pairs show the same pattern, while all 2096 other 2020 pairs have matching channels. Original CSV checks confirm these fields. A single timestamp without an independent transmission identifier cannot distinguish repetitions from metadata errors; all pairs remain evaluated.

July's 0.99 quantile interpolates between 2.120 and 3.471 ms; the 0.999 quantile lies between 2.136977 and 4.201604 s, explaining the jump to 3.308 s. Channel disagreement further limits the payload reference's physical meaning, precluding interpretation as clock drift alone.

A reference-free training lag histogram uses time, channel, spreading factor, and gateway ID with fixed support and peak-separation rules. Refinement contains 251 candidates with median 0.999985 s, matching the supervised offset numerically. Both fits give identical partitions on 23,677 reports, with TP 1411, FP 1, FN 6, and F1 0.9975. Population selection, scoring, and gate calibration remain reference-assisted (Supplementary Section S15).

The 2019-only fit has one positive pair from gateways absent in the 2020 tests, and April has no cross-gateway reference pair. Zero F1 combines sparse supervision and unsupported calibration rather than isolating drift. These retrospective checks diagnose preprocessing and reference construction.

\FloatBarrier
